\documentclass{article}
\usepackage{listings}
\usepackage{amsmath,amsthm}
\usepackage{amssymb}
%\PassOptionsToPackage{hyphens}{url}
\usepackage{hyperref}
\usepackage{xurl}


\title{Group projects Financial Modeling I}
\author{}
\date{}

\begin{document}
\maketitle

\textbf{Deadline: 30/04/2025, 23:59,  please submit your *.py or *.ipynb files and a pdf file with the answers 
       using the designated folder on the K2 site of the course.}



%\section*{Project 1: Team 5 (PROSPERI Armand,ORSAZ Léna,GRIFFOUL Victorine), Team 1 (MICHEL Bastien,ORBAN James,MENDY Gélaze, POIRIER Thibault)}

% \subsection*{Linear regression for independetly identically sampled data}
%%  The task is to  predict the logarithm of the hourly wages (in USD) (variable $lwage$) based on education (variable $educ$), work experience (variable $exper$), tenure (variable $tenure$) and gender (variable $female$). The gender is encoded in the variable $female$ which is $1$ for   females, zero for males. 
%% The .csv file is called \lsi
%%
%% \url{https://kaggle.com/datasets/02f69ffd32ee24c19c78526e7dee71cfe75a475a78c48b66c63fc11dc1d09e75}. 
%%
%% Linear regression for predicting the logarithm of wage. 
%% Use the independent variables $educ$, $exper$, $tenure$ and $female$
%% and use them to build models for predicting the logarithm of the wage ($lwage$).
%%
%% Reserve at least 20\% of your data for validation.
%%
%%  \begin{itemize}
%%  \item Choose randomly two sets of 100 data points from your data set, and for each data set compute the mean, variance, and the histogram
%%	of each of the variables  $educ$, $exper$, $tenure$,  $female$  and $lwage$. 
%%	Are these values similar for the two data sets ?  Can we say that each of these variables are independetly identically sampled ? 
%%  \item Check if there is any linear relationship between variables. 		  
%%  \item Calculate the $R^2$ score on the validation data set. Comment on the quality of the model.
%%  \item Evalute the confidence intervals for the estimated regression parameters. Are all regression coefficients statistically significant ? Explain your answer. 
%%  \item Interpret the regression coefficient for the variable $female$. What does it say about the influence of gender on wages ?
%%%  \item Try to estimate by linear regression the logarithm of the hourly wage using only gender. Comment on the result
%%%  \item In addition to the variables $educ$, $exper$,$female$, use the variables $female \cdot (educ - m)$ instead, where $m$ is the average number of formal education,m
%%
%%  \item  Try to use  square $exper$ (variable $expersq$) in addition to $educ$, $exper$, $tenure$ and $female$ to predict $lwage$. What do you see ? 
%%  \item  Use  first $\log(educ+1)$ and $\log(exper+1)$ and $\log((exper+1)*(educ+1))$ to predict $lwage$ (without constant coefficient, i.e. $\beta_0=0$), and then use only 
%%	  $\log(educ+1)$ and $\log(exper+1)$ to predict   $lwage$ (again without constant coefficient, i.e. $\beta_0=0$),.  For both models compute the confidence intervals of regression coefficients.
%%	   What can you conclude about the effect of $\log(exper+1)$ on $lwage$
%%	  from the first and from the second model ?   Can you explain it ?
%%	 Is there a  linear relationship between the various independent variables ? 
%%  \end{itemize}
%%
%% \subsection*{Time-series prediction}
%%   Use the data set  
%%
%%    \url{https://kaggle.com/datasets/c0375f8666f329bf473a90c56739d47f9c9351706018cebe96fc4c385ee91b1d} 
%%
%%    to predict the weekly 
%%   price index of the New York Stock Exchange using the price indices of the previous two weeks. The prince index is represented by
%%   variable $price$. Use linear regression. Consider detrending using affine dependence on time.  
%%   Evaluate the $R^2$ score and try taking into account several past values: price index of
%%   the previous week, previous two weeks, previous three weeks, etc. 


\section*{Project 1:  Group 6 (Maxime Yilmaz, Alexandre Paquet, Christian Cheng), Group 7 (Zuer Li,Qiuyan Li,LUONG Thi Lan Hue), Group 5 (Célia Dompiétrini, Loane Mboutou,  Zaineb Er-raji) }


 \subsection*{Linear regression for independetly identically sampled data}
  The task is to predict the logarithm of the house prices (variable $lprice$) using the following variables:
  logarithm of the lot size (variable $llotisize$), logarithm of the surface of the house (variabl $lsqrtft$), number of bedrooms (variable $bdrms$). 	
   The .csv file is called \lstinline$hprice1_merged.csv$.

   %\url{https://kaggle.com/datasets/eb36254d39a218a99ac169cedb764d10b6468df9aff710c1e3508651170ff355}. % \url{https://www.kaggle.com/datasets/mihalypetreczky/hpric1}.  

  \begin{enumerate}
   \item Choose randomly two sets of $40$ data points from your data set,  
	   and for each data set compute the mean, variance, and the histogram
	of each of the variables  $lprice$, $llotsize$, $lsqrtft$,  $bdrms$. 
	Are these values similar for the two data sets ?  Can we say that each of these variables are independetly identically sampled ? 
	Consider using hypothesis testing to compare the distribution of these two data sets.

  \item Estimate the coefficients of the linear regression from data. Reserve at least 20\% of your data for validation.
  \item Calculate the $R^2$ score and the average mean squared error. Comment on the quality of the model.
  \item Evalute the confidence intervals for the estimated regression parameters. Are all regression coefficients statistically significant ? Explain your answer. 
  \item Add the variable $\log(lotsize^2*sqrft)$ as a new variable, and construct a linear regression model predicting $lprice$ using
	 $llotsize$ and $lsqrtft$ and $bdrms$ and $\log(lotsize^2*sqrft)$. What can you say about the relevance of  $llotsize$ and $lsqrtft$ and $bdrms$
		  for predicting $lprice$ ? Is it consistent with the results of the previous model (the one without $\log(lotsize^2*sqrft)$).
		  Explain your answer. Is there a linear relationship among those variables ? 

  \end{enumerate}
 
 
 \subsection*{Time-series prediction}
 Use the data set  \lstinline$fertil31.csv$

   %\url{https://kaggle.com/datasets/045ebee9dbafdf9e30b88dcbcd64d33a29ae7a425e2e9d6b73cb9d0051569713} 
   to predict fertility rate (variable $gfr$)
    using the fertility rate of the previous years. 
    Use linear regression. Consider detrending using quadratic dependence on time.  
   Evaluate the $R^2$ score and try taking into account several past values: fertility rate of the previous year, of the previous two years, etc. 



\section*{Project 3: Group 8 (Killian Bruni,Maïlys Texier), Group 9 (ZHANG Ruge,CAO Ranjie) }
 \subsection*{Linear regression for independetly identically sampled data}
    The task is to predict house prices in different communities based on air pollution level (represented by the concentration of nitrogine oxid), 
distance from employment centers, average student-teacher ratio 
 in local schools and average number of rooms, and crime rate The variables are as follows:
    \begin{itemize}
    \item{$price$ \\}       
	median housing price
    \item{$crime$ \\}                    
		crimes committed per capita
    \item{$nox$ \\}                      
	nitrogine oxide, parts per 100 mill.
    \item{$rooms$ \\}                    
    avg number of rooms per house
    \item{$dist$ \\}                     
      weighted dist. to 5 employ centers
    \item{$stratio$ \\}                  average student-teacher ratio
    \item{$lprice$ \\}                   $\log(price)$
    \item{$lnox$ \\}                     $\log(nox)$
   \end{itemize}

   The .csv file is \lstinline$hprice2_mergede.csv$.

    %\url{https://kaggle.com/datasets/f06415afa98ac5a0a50ad7fdec3c4b526023dcf2a1b39435e6f2b6a43eddf88e}.% \url{https://www.kaggle.com/datasets/mihalypetreczky/hprice2}. 
    
 \begin{itemize}
	 \item Choose randomly two sets of 100 data points from your data set, and for each data set compute the mean, variance, and the histogram
	of each of the variables  $lprice$, $lnox$, $rooms$,  $dist$. 
	Are these values similar for the two data sets ?  Can we say that each of these variables are independetly identically sampled ? 
	Consider using hypothesis testing to compare the distribution of these two data sets.

   \item Find a linear regression model which uses the variables $lnox$, $dist$, $rooms$, $stratio$  in order to predict the variable $lprice$.
     When estimating the coefficients of the linear regression, reserve at least 20\% of your data for validation.
  \item Calculate the $R^2$ score. Comment on the quality of the model. Are all the independent variables also linearly independent ?
  \item Evalute the confidence intervals for the estimated regression parameters. Are all regression coefficients statistically significant ? Explain your answer.
  \item Find a  linear regression model which uses the variables $lnox$, $dist$, $rooms$, $stratio$ and the new variable
	$\log(nox^3*e^{dist+1})$ to predict $lprice$. Compute the confidence intervals of the estimated regression parameters.
	What can you say about the significance of the variables $lnox$, $dist$, $rooms$, $stratio$ in predicting $lprice$.
	Is it different from the significance in the first model ? If so, why ? Is there a linear relationship among the variables ? 
  \end{itemize}

   
 \subsection*{Time-series prediction}
   Use the data set \lstinline$perturbed_priceindex.csv$
    %\url{https://kaggle.com/datasets/1b473a28fa790cd0ae128969ea13f5ce68bc06e906614c1ae21bf554c99cdce3} 
    to predict the consumer price index 
   based on its values in the previous years.    Use linear regression.
   Consider detrending using quadratic dependence on time.  
   Evaluate the $R^2$ score and try taking into account several past values: price index  of the previous month, of the previous two months, etc. 




\section*{Project 4: Group 1 (Anouk Tourtelier,Lana Szymoniak,Florine Rousteau), Group 3 (Vincent MALRANGE,Victor Bouffant,Léo Rioufrays--Laplace) }
 \subsection*{Linear regression for independetly identically sampled data}
   The task to predict logarithm of wages based on education, experience and IQ, hours worked. The variables are as follows:
   \begin{itemize}
   \item{$lwage$ \\}                    logarithm monthly earnings
   \item{$hours$ \\}                    average weekly hours
   \item{$IQ$ \\}                      IQ score
   \item{$educ$ \\}                     years of education
   \item{$exper$ \\}                    years of work experience
   \end{itemize}


   The .csv file is \lstinline$Wage2_merged.csv$.

%   \url{https://kaggle.com/datasets/053a610ff5d8f491e3ee24c7eadc563cdd094ff395e72ba6734ebeec4ab1ca59}. %\url{https://www.kaggle.com/datasets/mihalypetreczky/wage2dataset}. 

  \begin{enumerate}
	 \item Choose randomly two sets of 100 data points from your data set, and for each data set compute the mean, variance, and the histogram
		 of each of the variables  $lwage$, $hours$, $IQ$,  $educ$, $exper$. 
	Are these values similar for the two data sets ?  Can we say that each of these variables are independetly identically sampled ? 
	Consider using hypothesis testing to compare the distribution of these two data sets.



\item Find a linear regression model for predicting $lwage$ using the variables $\log(hours+1)$, $\log(exper+1)$ and $educ$. 
   Leave at least 20\% of the data for validation. 
  \item Calculate the $R^2$ score. Comment on the quality of the model.
  \item Evalute the confidence intervals for the estimated regression parameters. Are all regression coefficients statistically significant ? Explain your answer.
  \item Build a linear regression model  for predicting $lwage$ using the variables $\log(hours+1)$, $\log(exper+1)$ and $educ$ and $IQ$
	  and $\log((exper+1)*\sqrt{(hours+1)})$. Compute the confidence intervals of the estimated regression parameters.
	What can you say about the significance of the variables $\log(hours+1)$, $\log(exper+1)$ and $educ$ and 
	in predicting $lwage$.
	Is it different from the significance in the first model ? If so, why ? Is there a linear relationship among the variables ? 

 \end{enumerate}


   
 \subsection*{Time-series prediction}
    Use the data set from \lstinline$pcip_filtered.csv$ to predict the percentage change in monthly
    industrial production with respect to the previous year (variable $pcip$). Use linear regression.
   Consider detrending using affine  dependence on time.  
   Evaluate the $R^2$ score and try taking into account several past values: the value of $pcip$ the past year, the year before, etc. 

\section*{Project 4: Group 2 (Thelma GRECO,Angelick ORANGE), Group  5 (Yipan Liu,Ruixuan Li,Jingyan Wu)} 
)}

 \subsection*{Linear regression for independetly identically sampled data}
    The task is to predict the percentage of students receiving a passing grade in math (variable $math10$) in function of 
    the following variables:
    \begin{itemize}
    \item{$lnchprg$}    
	    percentage of students getting lunch at school
     \item{$staff$}     
	    number of teachers per 1000 students
     \item{$expend$}    
	    expenditure per student
    \item{$salary$}
	    average teacher salary
    \item{$graduation$}
	    graduation rate (percentage)
    \item{$benefits$}
	    average benefits (outside salary) of teachers.
    \end{itemize}		    
     The .csv file is \lstinline$meap93.csv$.

     %\url{https://kaggle.com/datasets/c490c3b6eacc1c99cb84a5890c8646f82dc05ffe33a52c20b35c64b61cf730ef}. 

  \begin{enumerate}
	 \item Choose randomly two sets of 100 data points from your data set, and for each data set compute the mean, variance, and the histogram
		 of each of the variables $lnchprg$, $staff$, $expend$, $salary$, $math10$. 
	Are these values similar for the two data sets ?  Can we say that each of these variables are independetly identically sampled ? 
	Consider using hypothesis testing to compare the distribution of these two data sets.



\item Find a linear regression model for predicting $math10$ using the variables $lnchprg$, $\log(expend)$, $\log(salary)$. Find linear
	regression models for predicting $math10$ using only  $lnchprg$, only $\log(expend)$ and only $\log(salary)$. 
   Leave at least 20\% of the data for validation. 
  \item Calculate the $R^2$ score for each of the $4$ models. Comment on the quality of the model.
  \item Evalute the confidence intervals for the estimated regression parameters for each of the for $4$ models. Are all regression coefficients statistically significant ? Explain your answer.
   \item
	What can you say about the significance of the variables $\log(expend)$, $\log(salary)$ and $lnchprg$
	in predicting $math10$.
	Do the coefficients of $\log(expend)$, $lnchprg$, $\log(salary)$ change a lot when you use all $3$ variables to predict $math10$ 
	compared to the case when you use only one of them ? For example, is the coefficient
		  of $\log(expend)$ in the model $\widehat{math10}=\beta_0+\beta_1 \log(expend)$  
		    is similar to the coefficient of $\log(expend)$ in the model   $\widehat{math10}=\beta_0+\beta_1 \log(expend)+\beta_2 \log(salary)+\beta_3 lnchprg$.
	If so, why ? Is there a linear relationship among the variables ? 

 \end{enumerate}


 \subsection*{Time-series prediction}
     Use the data set \lstinline$barium_import_filtered.csv$
      %\url{https://kaggle.com/datasets/a0ad6cf9f223e382a09fc5e8602ea450c118165eb5a6b351e7b2875bb3c2c554} 
      to predict the annual barium import
    (variable $chnimp$). Use linear regression.
   Consider detrending using affine  dependence on time.  
   Evaluate the $R^2$ score and try taking into account several past values: the value of $chnimp$ the past year, the year before, etc. 




\end{document}
